Multimodal Large Language Models have demonstrated remarkable capabilities in multi-modal understanding and generation. However, they persistently suffer from hallucination issues, generating text that contradicts visual inputs. The fundamental root cause of these hallucinations lies in the cross-modal granularity mismatch during autoregressive decoding: while MLLMs encode images as continuous holistic features, they are forced to generate text as discrete, often semantically fragmented sub-word tokens. This structural misalignment makes it mathematically intractable to ground individual, meaningless tokens to concrete visual entities. Consequently, the model is forced into an unverified decoding state, where the statistical inertia of language priors inevitably breeds hallucinations. Existing mitigation strategies struggle to address this efficiently, as they typically require computationally expensive model retraining or incur severe inference overhead through complex multi-round verification. Motivated by this insight, we propose CHORD, a training-free, plug-and-play decoding framework designed to forcefully re-anchor text generation to visual reality via Object-Resonant Decoding. Inspired by the physical resonance process, CHORD intercepts the active decoding frontier and dynamically evaluates the alignment between candidate text spans and external visual evidence. A candidate span's generation probability is amplified only if it semantically ``resonates'' with a concrete visual object, whereas ungrounded hallucinations---driven merely by linguistic inertia---are decisively penalized. Furthermore, to limit overhead while still catching overconfident errors, we use continuous Top-$1$ monitoring with full-$K$ backoff, exploiting KV-cache reuse in the external text encoder so that expensive full-frontier scoring runs only when the top hypothesis fails a resonance check. Extensive experiments across multiple hallucination-focused benchmarks demonstrate that CHORD substantially reduces object hallucinations. As a model-agnostic framework, it requires no additional training and establishes a highly scalable paradigm for trustworthy MLLM deployment.
